%%%%%%%%%%%%%%%%%
% Nokia-targeted CV — based on altacv.cls
% Target: FW Development Engineer – L0 (Optical Networks)
%%%%%%%%%%%%%%%%

\PassOptionsToPackage{dvipsnames}{xcolor}

\documentclass[10pt,a4paper,ragged2e,withhyper]{altacv}

\newenvironment{sloppypar*}{\sloppy\ignorespaces}{\par}

\geometry{left=1.2cm,right=1.2cm,top=1cm,bottom=1cm,columnsep=0.75cm}

\usepackage{paracol}

\ifxetexorluatex
  \setmainfont{Roboto Slab}
  \setsansfont{Lato}
  \renewcommand{\familydefault}{\sfdefault}
\else
  \usepackage[rm]{roboto}
  \usepackage[defaultsans]{lato}
  \renewcommand{\familydefault}{\sfdefault}
\fi

\ifdarkmode%
    \definecolor{PrimaryColor}{HTML}{0A2540}
    \definecolor{SecondaryColor}{HTML}{1F4E79}
    \definecolor{ThirdColor}{HTML}{FF7A00}
    \definecolor{BodyColor}{HTML}{333333}
    \definecolor{EmphasisColor}{HTML}{111111}
    \definecolor{BackgroundColor}{HTML}{FFFFFF}
\else%
  \definecolor{PrimaryColor}{HTML}{111111}
  \definecolor{SecondaryColor}{HTML}{2563EB}
  \definecolor{ThirdColor}{HTML}{111111}
  \definecolor{BodyColor}{HTML}{555555}
  \definecolor{EmphasisColor}{HTML}{111111}
  \definecolor{BackgroundColor}{HTML}{FFFFFF}
\fi%

\colorlet{name}{PrimaryColor}
\colorlet{tagline}{SecondaryColor}
\colorlet{heading}{PrimaryColor}
\colorlet{headingrule}{ThirdColor}
\colorlet{subheading}{SecondaryColor}
\colorlet{accent}{SecondaryColor}
\colorlet{emphasis}{EmphasisColor}
\colorlet{body}{BodyColor}
\pagecolor{BackgroundColor}

\renewcommand{\namefont}{\Huge\rmfamily\bfseries}
\renewcommand{\personalinfofont}{\small\bfseries}
\renewcommand{\cvsectionfont}{\Large\bfseries}
\renewcommand{\cvsubsectionfont}{\large\bfseries}

\renewcommand{\itemmarker}{{\small\textbullet}}
\renewcommand{\ratingmarker}{\faCircle}

\begin{document}
    \name{Martim Silva}
    \tagline{Embedded Firmware Engineering Student | RTOS, Bare-Metal C/C++ \\ Peripheral Drivers and Real Time Systems}

    \photoL{4cm}{AI_good_face}

    \personalinfo{
        \email{Martim.Alexandre.Dras@gmail.com}\smallskip
        \phone{+351-936-495-912}
        \location{Lisbon, Portugal}\\
        \linkedin{martim-silva-tech-dev/}
        \github{XanderDras}
    }

    \makecvheader

    \columnratio{0.22}

    \begin{paracol}{2}
        % ----- TECH STACK -----
        \cvsection{SKILL STACK}
            \begin{sloppypar*}
                \cvtags{C/true, C++/true, FreeRTOS/true, Embedded Systems/true, I2C/true, SPI/true, UART/true, GPIO/true, Python, ESP-IDF, Git, GDB, Linux, Docker}
                \medskip
            \end{sloppypar*}
        % ----- TECH STACK -----

        % ----- LANGUAGES -----
        \cvsection{Languages}
            \cvlang{English}{C2}\\
            \divider

            \cvlang{Portuguese}{Native}
            \bigskip
        % ----- LANGUAGES -----

        % ----- EDUCATION -----
        \cvsection{Education}

        \textbf{Bachelor of Science in Informatics, Networks and Telecommunications Engineering}\\
        \smallskip
        Instituto Superior de Engenharia de Lisboa (ISEL)\\
        \smallskip
        Sep 2023 -- Jul 2026

        \divider

        \textbf{Vocational Training in Information Systems Management and Programming}\\
        \smallskip
        Escola Secundária Gago Coutinho\\
        \smallskip
        2020 -- 2023\\
        Student Excellence Award – Municipality of Vila Franca de Xira
        % ----- EDUCATION -----

        % ----- REFERENCES -----
        \cvsection{References}
            \cvref{Prof.\ Rui Duarte}{ISEL}{}
            \divider

            \cvref{Boss\ Paulo Gouveia}{YEPKIT}{}
        % ----- REFERENCES -----

        \newpage

        \switchcolumn

        % ----- ABOUT ME -----
        \cvsection{About Me}
            \begin{quote}
            Engineering student experienced in developing firmware for embedded systems using FreeRTOS and bare-metal C/C++. Worked on designing hardware peripheral drivers (I2C, SPI, UART, GPIO), scheduling real time tasks, and developing power management for a PocketQube satellite platform.
            \end{quote}
        % ----- ABOUT ME -----

        % ----- PROJECTS -----
        \cvsection{Engineering Projects}
            \cvevent{PocketQube -- Modular COTS Satellite Platform}{Academic Engineering Project}{Mar 2026 -- Jul 2026 (ONGOING)}{ISEL}
            \begin{itemize}
                \item Developing FreeRTOS based firmware in C for an embedded SoC; implementing real time tasks scheduling, interrupt handling, and acquiring telemetry and downlink over 433~MHz RF communication.
                \item Designing hardware peripheral drivers over I2C, SPI, UART, and GPIO for controlling sensors and RF modules; taking care of the power subsystem and building firmware with ESP-IDF cross compilation.
            \end{itemize}

            \cvevent{Smart Car -- Web-Controlled IoT Platform}{Academic Project}{Jan 2023 -- Jul 2023}{}
            \begin{itemize}
                \item Working on interfacing the GPIO of Raspberry Pi to control DC motors and acquire sensors data; implementing real time video streaming using OpenCV and WebSocket API.
            \end{itemize}
        % ----- PROJECTS -----

        % ----- EXPERIENCE -----
        \cvsection{Experience}
            \cvevent{Software Developer}{YEPKIT}{May 2023 -- Jul 2023}{Portugal (Hybrid)}
            \begin{itemize}
                \item Developing and containerising back end services in Node.js with Docker; applying modular approach to the development with Git workflow.
            \end{itemize}
            \divider

            \cvevent{Team Leader | CubeSat Portugal Bootcamp}{ISEL-Space Club}{April 2025}{Lisbon, Portugal}
            \begin{itemize}
                \item Leading a team of 30 individuals with embedded systems, RF communications and mission planning responsibilities, to build a prototype CubeSat under supervision of professionals of Portuguese Space Agency.
            \end{itemize}

            \cvevent{Team Leader | Robotics Competitions}{Escola Secundária Gago Coutinho}{2022 -- 2023}{Portugal}
            \begin{itemize}
                \item Represented Portugal at International First Global 2022 (Geneva), European RoboCup 2022, Robótica 2023, and CanSat 2022.
                \item Designing and programming robots in C/C++ for autonomous operation on Arduino microcontrollers.
            \end{itemize}
        % ----- EXPERIENCE -----

    \end{paracol}
\end{document}
